\chapter[Functional Framework]{Functional Framework for the UOD}

The UOD admits a natural functional interpretation on the space of
log-likelihood fields.  For fixed $P$ and $Q$, define the
log-likelihood ratio $a_i = \log(P_i/Q_i)$ and consider it as a random
variable with respect to the uniform reference $U_n$.  The UOD is then
the variance of this random variable:

\[
\mathcal D(P,Q) = \operatorname{Var}_{U_n}(a) = \frac1n\sum_i a_i^2 - \Bigl(\frac1n\sum_i a_i\Bigr)^2.
\]

This variance interpretation unifies three perspectives: the first
moment (the KL divergence $D_{\mathrm{KL}}(P\|U_n) = \mathbb E_{U_n}[a]$),
the reverse KL, and the second centered moment (the UOD).  All classical
information measures can be expressed as functions of these moments
evaluated at different reference measures $Q$.

The functionals studied in this book---Shannon, R\'enyi, Tsallis,
min-entropy, Bayes vulnerability, guessing entropy, KL, $\chi^2$,
Hellinger, $\alpha$-divergences, $\beta$-divergences---are each a
special case of the general framework:

\[
S(P) = \Phi\!\left(\sum_i f(P_i, Q_i)\right),
\]

where $f$ is a smooth kernel and $\Phi$ is a normalising function.
The WIS geometry controls the gradient and Hessian of $S$ through
the chain rule, and the universal bounds of Chapter~14 apply
automatically once $f$ and $\Phi$ are twice differentiable.